\documentclass{ctexart}
\usepackage{amsmath,amssymb}
\begin{document}
\section*{卡车与配送机器人协同路径：完整操作模型}
以原论文第4--5页式(1)--(10)纠正旧02的抽取错误。节点集合
$V=\{v_0\}\cup C$，$v_0$为仓库，客户各有一个包裹。卡车和一个机器人
从仓库共同出发并返回。坐标单位km，速度km/h，所有作业时间单位秒；
$d^T_{ab}=3600D_{ab}/v_T$，$d^R_{ab}=3600D_{ab}/v_R$。
输入的完整距离矩阵由既有坐标的欧氏距离得到。客户服务时间为$t_s$，
仓库没有客户服务。装卸机器人分别耗时$t_l,t_u$，机器人最多携带$p$件。

\subsection*{可行操作及费用}
一个操作$o=(s_o,H_o,J_o,e_o)$包含共同出发、重会节点，
有序卡车内部客户序列$H_o$和有序机器人客户序列$J_o$；
内部客户互不重复、互不重叠，且$|J_o|\le p$。
机器人为空时随车行驶。操作集合$O$包含所有合法此类操作，
不限于算法实际枚举的前向操作。
令$P^T_o=(s_o,H_o,e_o)$、$P^R_o=(s_o,J_o,e_o)$，有
\[
 T_o=\sum_{(a,b)\in P^T_o}d^T_{ab}
       +t_s(|H_o|+\mathbf1_{\{e_o\ne v_0\}}),\qquad
 R_o=\begin{cases}
 \sum_{(a,b)\in P^R_o}d^R_{ab}+|J_o|t_s+t_l+t_u,&J_o\ne\varnothing,\\
 0,&J_o=\varnothing,
 \end{cases}\quad c_o=\max(T_o,R_o).
\]
这里求和按路径的相邻弧进行；在终点客户服务记入前一操作，
以免连续操作重复计费。仓库无服务是数据边界约定。

\subsection*{决策与约束}
$x_o$选择操作，$y_v$表示作为操作分界的共同节点。
$O(v)$包含访问$v$的操作；$O^-(v)$以$v$出发，
$O^+(v)$以$v$结束。$O^+(S)$为从$V\setminus S$进入$S$的操作。
\begin{align*}
\min\quad &\sum_{o\in O}c_ox_o\\
\text{s.t.}\quad
&\sum_{o\in O(v)}x_o\ge1 &&v\in V,\\
&\sum_{o\in O^-(v)}x_o\le y_v &&v\in V,\\
&\sum_{o\in O^-(v)}x_o=\sum_{o\in O^+(v)}x_o &&v\in V,\\
&\sum_{o\in O^+(S)}x_o\ge y_v
 &&\varnothing\ne S\subseteq V\setminus\{v_0\},\ v\in S,\\
&\sum_{o\in O^+(v_0)}x_o\ge1,\qquad y_{v_0}=1,\\
&x_o,y_v\in\{0,1\}.
\end{align*}
旧02的覆盖$\le1$、连通割右端$\sum_{v\in S}y_v$以及若干等号
与源文不一致，按源文分别纠正。客户在交付候选中恰好服务一次，
因此也满足上述至少一次覆盖。

\subsection*{算法与证书}
代码先构造包含全部客户的卡车闭合路径，再做2-opt和局部操作动态规划。
最多跨7个路径位置、机器人长序列使用近邻顺序，均为搜索策略；
这些条件用于生成候选操作序列。
交付的有序操作构成单一闭合链，节点无重复分界，因此满足所有连通割，
无需枚举指数多个集合。另行检查客户覆盖、容量、费用和首尾衔接。
每实例实际3600秒预算；返回可行上界，无最优性证明。
\end{document}
